
\providecommand{\GL}{\operatorname{GL}}
\providecommand{\PGL}{\operatorname{PGL}}
\providecommand{\AGL}{\operatorname{AGL}}
\providecommand{\Hol}{\operatorname{Hol}}
\providecommand{\Sym}{\operatorname{Sym}}

\chapter{Linear and Affine Permutation Groups}

\section{Multiple Transitivity}

\begin{definition}[\(k\)-transitive action]
Let \(G\leq \Sym(\Omega)\), and let \(k\geq 1\). We say that \(G\) is
\emph{\(k\)-transitive} on \(\Omega\) if for any two ordered \(k\)-tuples
\[
    (\omega_1,\ldots,\omega_k),\qquad (\omega_1',\ldots,\omega_k')
\]
of pairwise distinct points of \(\Omega\), there exists \(g\in G\) such that
\[
    \omega_i^g=\omega_i'
    \qquad\text{for all }1\leq i\leq k.
\]
Equivalently,
\[
    (\omega_1,\ldots,\omega_k)^g
    =
    (\omega_1',\ldots,\omega_k').
\]
\end{definition}

\begin{definition}[Sharply \(k\)-transitive action]
Let \(G\leq \Sym(\Omega)\). We say that \(G\) is
\emph{sharply \(k\)-transitive} on \(\Omega\) if for any two ordered \(k\)-tuples
of pairwise distinct points
\[
    (\omega_1,\ldots,\omega_k),\qquad
    (\omega_1',\ldots,\omega_k'),
\]
there exists \emph{exactly one} element \(g\in G\) such that
\[
    \omega_i^g=\omega_i'
    \qquad\text{for all }1\leq i\leq k.
\]
\end{definition}

\begin{proposition}[Stabilizer criterion for \(k\)-transitivity]
Let \(G\leq \Sym(\Omega)\), let \(\omega\in \Omega\), and let \(k\geq 2\).
Then \(G\) is \(k\)-transitive on \(\Omega\) if and only if
\[
    G \text{ is transitive on } \Omega
    \quad\text{and}\quad
    G_\omega \text{ is }(k-1)\text{-transitive on }\Omega\setminus\{\omega\}.
\]
\end{proposition}

\begin{proof}
Suppose first that \(G\) is \(k\)-transitive. Then \(G\) is certainly transitive.
Moreover, if
\[
    (\alpha_2,\ldots,\alpha_k),\qquad
    (\beta_2,\ldots,\beta_k)
\]
are two ordered \((k-1)\)-tuples of pairwise distinct points in
\(\Omega\setminus\{\omega\}\), then \(k\)-transitivity gives some \(g\in G\) such that
\[
    (\omega,\alpha_2,\ldots,\alpha_k)^g
    =
    (\omega,\beta_2,\ldots,\beta_k).
\]
In particular, \(\omega^g=\omega\), so \(g\in G_\omega\). Hence \(G_\omega\)
is \((k-1)\)-transitive on \(\Omega\setminus\{\omega\}\).

Conversely, suppose that \(G\) is transitive on \(\Omega\) and \(G_\omega\) is
\((k-1)\)-transitive on \(\Omega\setminus\{\omega\}\). Let
\[
    (\omega_1,\ldots,\omega_k),\qquad
    (\omega_1',\ldots,\omega_k')
\]
be ordered \(k\)-tuples of pairwise distinct points. Since \(G\) is transitive, choose
\(g_1,g_2\in G\) such that
\[
    \omega_1^{g_1}=\omega,\qquad
    (\omega_1')^{g_2}=\omega.
\]
Then
\[
    \omega_2^{g_1},\ldots,\omega_k^{g_1},
    \qquad
    (\omega_2')^{g_2},\ldots,(\omega_k')^{g_2}
\]
are two ordered \((k-1)\)-tuples of pairwise distinct points in
\(\Omega\setminus\{\omega\}\). By the hypothesis on \(G_\omega\), there exists
\(h\in G_\omega\) such that
\[
    \bigl(\omega_i^{g_1}\bigr)^h=(\omega_i')^{g_2}
    \qquad\text{for }2\leq i\leq k.
\]
Since \(h\) fixes \(\omega\), we also have
\[
    \bigl(\omega_1^{g_1}\bigr)^h=\omega=(\omega_1')^{g_2}.
\]
Therefore
\[
    (\omega_1,\ldots,\omega_k)^{g_1 h g_2^{-1}}
    =
    (\omega_1',\ldots,\omega_k'),
\]
so \(G\) is \(k\)-transitive.
\end{proof}

\begin{definition}[Rank of a transitive permutation group]
Let \(G\leq \Sym(\Omega)\) be transitive. The \emph{rank} of \(G\) on \(\Omega\)
is the number of \(G_\omega\)-orbits on \(\Omega\), where \(\omega\in \Omega\).
Equivalently, it is the number of \(G\)-orbits on \(\Omega\times\Omega\).
\end{definition}

\begin{remark}
A transitive group \(G\leq \Sym(\Omega)\) is \(2\)-transitive if and only if it has
rank \(2\). Indeed, one orbit on \(\Omega\times\Omega\) is the diagonal
\[
    \{(\omega,\omega)\mid \omega\in \Omega\},
\]
and \(2\)-transitivity is precisely the assertion that all off-diagonal pairs lie in
one orbit.
\end{remark}

\begin{lemma}[Counting criterion for sharp transitivity]
Assume \(\Omega\) is finite. If \(G\) is \(k\)-transitive on \(\Omega\), then \(G\) is
sharply \(k\)-transitive if and only if
\[
    |G|=|\Omega|(|\Omega|-1)\cdots(|\Omega|-k+1).
\]
\end{lemma}

\begin{proof}
Let \(\Omega^{(k)}\) denote the set of ordered \(k\)-tuples of pairwise distinct
points of \(\Omega\). Then
\[
    |\Omega^{(k)}|
    =
    |\Omega|(|\Omega|-1)\cdots(|\Omega|-k+1).
\]
If \(G\) is \(k\)-transitive, then \(G\) acts transitively on \(\Omega^{(k)}\).
Thus, for a fixed tuple \(\mathbf{\omega}\in \Omega^{(k)}\),
\[
    |\Omega^{(k)}|=[G:G_{\mathbf{\omega}}].
\]
The action is sharply \(k\)-transitive exactly when
\(G_{\mathbf{\omega}}=1\), which is equivalent to
\[
    |G|=|\Omega^{(k)}|.
\]
\end{proof}

%----------------------------------------------------------------------------------------
\section{The Natural Action of \texorpdfstring{$\GL(V)$}{GL(V)}}

Let \(\F=\F_q\) be a finite field, let \(V=\F^d\) with \(d>1\), and let
\[
    G=\GL(V).
\]
The group \(G\) acts naturally on \(V\). There are two orbits:
\[
    \{0\}
    \qquad\text{and}\qquad
    \Omega:=V^\#:=V\setminus\{0\}.
\]

\begin{proposition}
The action of \(\GL(V)\) on \(\Omega=V^\#\) is transitive.
\end{proposition}

\begin{proof}
Let \(u,v\in V^\#\). Extend \(u\) to a basis
\[
    u,u_2,\ldots,u_d
\]
of \(V\), and extend \(v\) to a basis
\[
    v,v_2,\ldots,v_d
\]
of \(V\). There is a unique linear isomorphism \(g\in \GL(V)\) such that
\[
    u^g=v,\qquad u_i^g=v_i\quad(2\leq i\leq d).
\]
Hence \(u\) and \(v\) lie in the same \(G\)-orbit.
\end{proof}

\begin{proposition}
If \(q>2\), then the action of \(\GL(V)\) on \(\Omega=V^\#\) is imprimitive.
\end{proposition}

\begin{proof}
For \(v\in V^\#\), define
\[
    B_v:=\F^\times v=\{\lambda v\mid \lambda\in \F^\times\}.
\]
Let
\[
    \mathcal B=\{B_v\mid v\in V^\#\}.
\]
Then \(\mathcal B\) is the set of nonzero points on the one-dimensional subspaces
of \(V\). Since \(q>2\),
\[
    |B_v|=q-1>1.
\]
Also, since \(d>1\),
\[
    |B_v|=q-1<q^d-1=|\Omega|.
\]
Thus the blocks are nontrivial and proper.

For any \(g\in \GL(V)\),
\[
    B_v^g
    =
    \{\lambda v^g\mid \lambda\in \F^\times\}
    =
    B_{v^g}.
\]
Hence \(g\) permutes the members of \(\mathcal B\). Therefore \(\mathcal B\) is a
nontrivial block system for the action of \(\GL(V)\) on \(\Omega\), and the action is
imprimitive.
\end{proof}

\begin{proposition}
If \(q=2\), then the action of \(\GL(V)\) on \(\Omega=V^\#\) is \(2\)-transitive.
\end{proposition}

\begin{proof}
Let
\[
    (u_1,u_2),\qquad (v_1,v_2)
\]
be ordered pairs of distinct points in \(\Omega\). Since \(q=2\), the only nonzero
scalar is \(1\). Hence \(u_1\neq u_2\) implies that \(u_1,u_2\) are linearly
independent. Similarly, \(v_1,v_2\) are linearly independent.

Extend
\[
    u_1,u_2
\]
to a basis
\[
    u_1,u_2,\ldots,u_d
\]
of \(V\), and extend
\[
    v_1,v_2
\]
to a basis
\[
    v_1,v_2,\ldots,v_d
\]
of \(V\). There exists \(g\in \GL(V)\) such that
\[
    u_i^g=v_i
    \qquad\text{for all }1\leq i\leq d.
\]
In particular,
\[
    (u_1,u_2)^g=(v_1,v_2).
\]
Thus \(\GL(V)\) is \(2\)-transitive on \(V^\#\).
\end{proof}

\begin{proposition}[Rank of \(\GL(V)\) on \(V^\#\)]
Assume \(q>2\). Then the rank of \(\GL(V)\) on \(\Omega=V^\#\) is \(q\).
More precisely, the \(G\)-orbits on \(\Omega\times\Omega\) are:
\[
    \Delta=\{(v,v)\mid v\in \Omega\},
\]
\[
    \mathcal O_\lambda
    =
    \{(v,\lambda v)\mid v\in \Omega\},
    \qquad \lambda\in \F^\times\setminus\{1\},
\]
and
\[
    \mathcal O_{\mathrm{ind}}
    =
    \{(u,v)\in \Omega\times\Omega\mid \langle u\rangle\neq \langle v\rangle\}.
\]
Therefore
\[
    \operatorname{rank}(G,\Omega)
    =
    1+(q-2)+1
    =
    q.
\]
\end{proposition}

\begin{proof}
The diagonal \(\Delta\) is clearly one orbit.

Now fix \(\lambda\in \F^\times\setminus\{1\}\). If \(u,v\in \Omega\), choose
\(g\in \GL(V)\) such that \(u^g=v\). Then
\[
    (u,\lambda u)^g=(v,\lambda v).
\]
Thus \(\mathcal O_\lambda\) is an orbit. Moreover, if
\[
    (u,\lambda u)^g=(v,\mu v),
\]
then
\[
    \lambda u^g=\mu v.
\]
Since \(u^g=v\), we get \(\lambda=\mu\). Hence the orbits
\(\mathcal O_\lambda\) are distinct as \(\lambda\) varies.

Finally, suppose \(\langle u\rangle\neq \langle v\rangle\) and
\(\langle u'\rangle\neq \langle v'\rangle\). Then \(u,v\) are linearly independent,
and \(u',v'\) are linearly independent. Extend both pairs to bases of \(V\), and
choose \(g\in \GL(V)\) mapping the first basis to the second. Then
\[
    (u,v)^g=(u',v').
\]
Thus all pairs with distinct one-dimensional spans form one orbit.

Hence the rank is
\[
    1+(q-2)+1=q.
\]
\end{proof}

%----------------------------------------------------------------------------------------
\section{The Projective Action}

Continue to let \(V=\F_q^d\) with \(d>1\), and let
\[
    G=\GL(V).
\]
Let
\[
    \mathcal B
    =
    \{B_v=\F^\times v\mid v\in V^\#\}.
\]
We identify \(\mathcal B\) with the set of one-dimensional subspaces of \(V\), i.e.
the points of the projective space associated with \(V\).

\begin{proposition}
The action of \(\GL(V)\) on \(\mathcal B\) is transitive, and the kernel of this
action is
\[
    Z(\GL(V))=\{\lambda I_V\mid \lambda\in \F^\times\}.
\]
Consequently the induced faithful permutation group is
\[
    \GL(V)/Z(\GL(V))=\PGL(V).
\]
\end{proposition}

\begin{proof}
Transitivity follows from the transitivity of \(\GL(V)\) on nonzero vectors:
if \(B_u,B_v\in \mathcal B\), choose \(g\in \GL(V)\) with \(u^g=v\). Then
\[
    B_u^g=B_v.
\]

Every scalar map \(\lambda I_V\) fixes every one-dimensional subspace, so
\[
    Z(\GL(V))\leq \Ker\bigl(\GL(V)\curvearrowright \mathcal B\bigr).
\]

Conversely, suppose \(g\in \GL(V)\) fixes every one-dimensional subspace.
Then for each \(0\neq v\in V\), there exists \(\lambda_v\in \F^\times\) such that
\[
    v^g=\lambda_v v.
\]
If \(u,v\) are linearly independent, then \(u+v\neq 0\), and \(g\) also fixes the
line \(\F(u+v)\). Hence
\[
    (u+v)^g=\lambda_{u+v}(u+v).
\]
But
\[
    (u+v)^g=u^g+v^g=\lambda_u u+\lambda_v v.
\]
Since \(u,v\) are linearly independent, we get
\[
    \lambda_u=\lambda_v=\lambda_{u+v}.
\]
It follows that the scalar \(\lambda_v\) is independent of \(v\), so \(g=\lambda I_V\).
Thus the kernel is exactly the scalar subgroup \(Z(\GL(V))\).
\end{proof}

\begin{proposition}
The group \(\PGL(V)\) is \(2\)-transitive on \(\mathcal B\).
\end{proposition}

\begin{proof}
Let
\[
    (B_{u_1},B_{u_2}),\qquad (B_{v_1},B_{v_2})
\]
be ordered pairs of distinct points of \(\mathcal B\). Since the projective points are
distinct, \(u_1,u_2\) are linearly independent, and \(v_1,v_2\) are linearly
independent.

Extend \(u_1,u_2\) to a basis of \(V\), and extend \(v_1,v_2\) to a basis of
\(V\). Then there exists \(g\in \GL(V)\) such that
\[
    u_1^g=v_1,\qquad u_2^g=v_2.
\]
Therefore
\[
    (B_{u_1},B_{u_2})^g=(B_{v_1},B_{v_2}).
\]
Passing to the quotient by the kernel, \(\PGL(V)\) is \(2\)-transitive on
\(\mathcal B\).
\end{proof}

\begin{theorem}
The group \(\PGL(V)\) is not \(3\)-transitive on \(\mathcal B\) if and only if
\[
    d>2.
\]
Equivalently, \(\PGL(V)\) is \(3\)-transitive on \(\mathcal B\) exactly when
\(d=2\).
\end{theorem}

\begin{proof}
First suppose \(d>2\). Choose linearly independent vectors
\[
    e_1,e_2,e_3\in V.
\]
Then the ordered triple
\[
    (\F e_1,\F e_2,\F e_3)
\]
has representatives that are linearly independent. On the other hand,
\[
    (\F e_1,\F e_2,\F(e_1+e_2))
\]
is an ordered triple of distinct projective points whose representatives are linearly
dependent. Linear maps preserve linear dependence and linear independence.
Therefore these two triples cannot lie in the same \(\PGL(V)\)-orbit. Thus
\(\PGL(V)\) is not \(3\)-transitive.

Conversely, suppose \(d=2\). Let
\[
    (\F u_1,\F u_2,\F u_3),\qquad
    (\F v_1,\F v_2,\F v_3)
\]
be ordered triples of distinct projective points. Since \(\F u_1\neq \F u_2\),
the vectors \(u_1,u_2\) form a basis of \(V\). Similarly, \(v_1,v_2\) form a
basis of \(V\).

Since \(\F u_3\) is distinct from both \(\F u_1\) and \(\F u_2\), we can write
\[
    u_3=a u_1+b u_2
\]
with \(a,b\in \F^\times\). Similarly,
\[
    v_3=c v_1+d v_2
\]
with \(c,d\in \F^\times\).

Define \(g\in \GL(V)\) by
\[
    u_1^g=\frac{c}{a}v_1,
    \qquad
    u_2^g=\frac{d}{b}v_2.
\]
Then
\[
    u_3^g
    =
    (a u_1+b u_2)^g
    =
    c v_1+d v_2
    =
    v_3.
\]
Hence
\[
    (\F u_1,\F u_2,\F u_3)^g
    =
    (\F v_1,\F v_2,\F v_3).
\]
Thus \(\PGL(V)\) is \(3\)-transitive when \(d=2\).
\end{proof}

\begin{remark}
When \(d=2\), the action of \(\PGL_2(q)\) on the projective line
\[
    \mathbb P^1(\F_q)
\]
is in fact sharply \(3\)-transitive. Indeed,
\[
    |\PGL_2(q)|=q(q^2-1),
\]
and the number of ordered triples of distinct points of \(\mathbb P^1(\F_q)\) is
\[
    (q+1)q(q-1)=q(q^2-1).
\]
\end{remark}

%----------------------------------------------------------------------------------------
\section{Affine Permutation Groups}

Let
\[
    N=(\Z/p\Z)^d
\]
where \(p\) is prime and \(d>1\). We regard \(N\) as the additive group of the
vector space \(\F_p^d\). Then
\[
    \Aut(N)=\GL_d(p).
\]
Set
\[
    H=\Aut(N)=\GL_d(p),
    \qquad
    G=N:H=\Hol(N).
\]
Thus
\[
    G=\F_p^d:\GL_d(p),
\]
the affine general linear group.

Let \(\Omega=N\). We define an action of \(G\) on \(\Omega\) as follows. If
\(x\in N\), \(y\in H\), and \(w\in N\), then
\[
    w^{xy}:=(w+x)^y=w^y+x^y.
\]
Here \(x\) acts as a translation and \(y\) acts as a linear transformation.

\begin{theorem}
The affine group
\[
    G=\F_p^d:\GL_d(p)
\]
is \(2\)-transitive on \(\Omega=\F_p^d\).
Furthermore, if \(p=2\), then \(G\) is \(3\)-transitive.
\end{theorem}

\begin{proof}
We first prove \(2\)-transitivity. Let
\[
    (u_1,u_2),\qquad (v_1,v_2)
\]
be ordered pairs of distinct points of \(N\). Then
\[
    u_2-u_1\neq 0,
    \qquad
    v_2-v_1\neq 0.
\]
Since \(\GL_d(p)\) is transitive on nonzero vectors, choose \(y\in \GL_d(p)\)
such that
\[
    (u_2-u_1)^y=v_2-v_1.
\]
Now choose
\[
    x=v_1^{y^{-1}}-u_1.
\]
Then
\[
    (u_1+x)^y
    =
    \bigl(u_1+v_1^{y^{-1}}-u_1\bigr)^y
    =
    v_1,
\]
and
\[
    (u_2+x)^y
    =
    \bigl(u_2-u_1+v_1^{y^{-1}}\bigr)^y
    =
    (u_2-u_1)^y+v_1
    =
    v_2.
\]
Thus
\[
    (u_1,u_2)^{xy}=(v_1,v_2),
\]
so \(G\) is \(2\)-transitive.

Now assume \(p=2\). The stabilizer of \(0\in N\) in \(G\) is
\[
    G_0=\GL_d(2).
\]
By the earlier result, \(\GL_d(2)\) is \(2\)-transitive on
\[
    N\setminus\{0\}.
\]
Since \(G\) is transitive on \(N\), the stabilizer criterion for multiple transitivity
implies that \(G\) is \(3\)-transitive on \(N\).
\end{proof}

\begin{remark}
The translation subgroup \(N\) is regular on \(\Omega=N\). It is also a minimal
normal subgroup of \(G\). Indeed, the normal subgroups of \(N\) that are normalized
by \(H=\GL_d(p)\) are precisely the \(H\)-invariant subspaces of \(\F_p^d\), and
\(\GL_d(p)\) acts irreducibly on the natural module. Hence the only such subgroups
are \(0\) and \(N\).

Moreover, every nontrivial normal subgroup of \(G\) contains \(N\). Therefore every
nontrivial normal subgroup of \(G\) is transitive on \(\Omega\). Thus \(G\) is a
quasiprimitive affine permutation group.
\end{remark}

%----------------------------------------------------------------------------------------
\section{A Sharply \texorpdfstring{$2$}{2}-Transitive Affine Example}

The full affine group \(\AGL_d(p)=\F_p^d:\GL_d(p)\) is often much larger than is
needed for \(2\)-transitivity. The following construction gives a smaller sharply
\(2\)-transitive group.

\begin{example}[The group \(\AGL_1(q)\)]
Let
\[
    F=\F_q,
    \qquad q=p^f.
\]
Then the additive group of \(F\) satisfies
\[
    F^+\cong (\Z/p\Z)^f,
\]
and the multiplicative group satisfies
\[
    F^\times\cong C_{q-1}.
\]
Choose a generator \(\varphi\) of \(F^\times\). Then
\[
    F^\times=\langle \varphi\rangle.
\]

The group \(F^\times\) acts on \(F^+\) by field multiplication:
\[
    x^\varphi=x\varphi
    \qquad (x\in F).
\]
This is an automorphism of the additive group \(F^+\). Hence
\[
    \langle \varphi\rangle\leq \GL_f(p),
\]
and we can form
\[
    G=F^+:\langle \varphi\rangle
    \leq
    \Hol(F^+)
    =
    F^+:\GL_f(p).
\]
Equivalently, \(G\) is the group of affine transformations
\[
    x\mapsto ax+b,
    \qquad
    a\in F^\times,\ b\in F.
\]
This group is usually denoted
\[
    \AGL_1(q).
\]

We claim that \(G\) is sharply \(2\)-transitive on \(F\).

First, \(F^\times\) is transitive on \(F^\times=F\setminus\{0\}\). Indeed, if
\(u,v\in F^\times\), then there exists \(a\in F^\times\) such that
\[
    ua=v.
\]
Therefore the stabilizer of \(0\) in \(G\) is transitive on \(F\setminus\{0\}\).
Since the translation subgroup \(F^+\) is transitive on \(F\), it follows that
\(G\) is \(2\)-transitive on \(F\).

Finally,
\[
    |G|=|F^+|\,|F^\times|=q(q-1).
\]
The number of ordered pairs \((x,y)\in F\times F\) with \(x\neq y\) is also
\[
    q(q-1).
\]
By the counting criterion, \(G\) is sharply \(2\)-transitive.

Equivalently, this uniqueness can be checked directly: given \(u_1\neq u_2\) and
\(v_1\neq v_2\) in \(F\), the affine map
\[
    x\mapsto ax+b
\]
sending \(u_1\mapsto v_1\) and \(u_2\mapsto v_2\) is uniquely determined by
\[
    a=\frac{v_2-v_1}{u_2-u_1},
    \qquad
    b=v_1-a u_1.
\]
\end{example}
